% SOURCE_AUTHORITY: sga5_fr_workpass.tex, lines 11518--11568 (LF-normalized unit).
% SOURCE_SHA256: 791F4EFFC5E02832D5D77ED03518C8156D6F07E4C8238B03545DB93D883FBB28
\subsection*{9.10}

Con las notaciones de (9.9), la aplicación
\[
\omega=\lambda\circ\lambda'
\]
es un proyector que verifica $\operatorname{Im}(\omega)=\operatorname{Im}(\lambda)$. Por consiguiente, $\mu$ induce un isomorfismo
\[
\bar\mu:\Ker(\omega)\xrightarrow{\sim}C^p(X').
\]
Pero
\[
\Ker(\omega)=\Ker(g_*)\oplus C^p(X)\hookrightarrow C^{p-1}(Y')\oplus C^p(X).
\tag{9.10.1}
\]
En efecto, un elemento $(x,y)$ de $C^{p-1}(Y')\oplus C^p(X)$ pertenece a $\Ker(\omega)$ si y solo si
\[
g^*g_*(x)\cdot c_{d-1}(\check F)=-i_*g_*(x)=0,
\]
es decir, teniendo en cuenta (9.1.1), $g_*(x)=0$. Por otra parte, la estructura del anillo de Chow de un fibrado proyectivo y, en particular, las relaciones
\[
g_*(\xi^i)=0\quad\text{para }i<d-1,\qquad g_*(\xi^{d-1})=1,
\]
muestran que
\[
\Ker(g_*)=g^*C^{p-1}(Y)\oplus g^*C^{p-2}(Y)\xi\oplus\cdots\oplus g^*C^{p-d+1}(Y)\xi^{d-2}.
\]
En otras palabras, $\bar\mu$ puede interpretarse como un isomorfismo
\[
C^{p-1}(Y)\oplus C^{p-2}(Y)\oplus\cdots\oplus C^{p-d+1}(Y)\oplus C^p(X)
\xrightarrow{\sim}C^p(X'),
\tag{9.10.2}
\]
cuyo isomorfismo inverso es fácil de explicitar.

Como me ha señalado Murre, el enunciado (9.10.2) fue demostrado para $d=2$ por Samuel, en \emph{Relations d'équivalence}, Proceedings of the International Congress, Edinburgh 1958.

Naturalmente, en el marco de la cohomología étale o $\ell$-ádica se tiene un enunciado análogo a (9.10.2), que esta vez se deduce de (8.5).

\subsection*{Bibliografía}

\begin{enumerate}[label={[\arabic*]}]
\item Grothendieck, \emph{La théorie des classes de Chern}, Bull. Soc. Math. de France \textbf{86} (1958), pp. 137--154.
\item Hodge, \emph{The theory and applications of harmonic integrals}, Cambridge University Press, 1952.
\item Hirzebruch, \emph{Topological Methods in Algebraic Geometry}, Springer Verlag.
\item Lefschetz, \emph{Algebraic Geometry}, Princeton University Press, 1953.
\item Sra. Raynaud, \emph{Modules projectifs universels}, Inventiones Mathematicae vol. 6. fasc. 1 (1968), pp. 1 a 26.
\item[{[TF]}] Godement, \emph{Théorie des faisceaux}.
\item[{[CH]}] \emph{Anneaux de Chow et applications}, Séminaire Chevalley 1958.
\item[{[6]}] \emph{Théorème de Lefschetz et critères de dégénérescence}, Pub. I.H.E.S. n$^\circ$ 35.
\end{enumerate}
